\scanpage{48}\mspage{47}

\deleted{\uline{Corollaire 2} Soient $F,\unsure{C|G}\in\mathcal F$, avec $G<F$. Alors}
\[
\deleted{\exists\,\unsure{,}\,x\in\mathrm{omb}(F),\ \text{t.q.}}
\]
\deleted{\unsure{J'arrive|J'en viens|Je reviens|J'arriverai} à la notion d'atelier \uline{ensembliste}} \deleted{\unsure{Déf}}

Je reprends les \unsure{sorites} des \deleted{homo}morphismes d'\unsure{Ateliers|ateliers}

\uline{Définition} Soient $\mathcal F',\deleted{\unsure{G|\mathcal G}}\mathcal F$ deux ateliers. On appelle \deleted{\unsure{prémo}} \uline{quasi-morphisme} \deleted{de $\mathcal F$} (en attendant une notion plus forte de morphisme !) de $\mathcal F'$ dans $\mathcal F$, une application
% slanted
\margin{il vaudrait peut-être mieux de parler de \uline{co}morphismes (\unsure{évt.|evt.} ``quasi''), cf plus bas.}
\[
f:\mathcal F'\longrightarrow\mathcal F
\]
qui est

a) croissante pour $\leq$ et $\ll$

b) commute aux Sup \deleted{\unsure{même} (\unsure{majorés|compris})} (et transforme \inserted{donc} $\varnothing_{\mathcal F'}$ en $\varnothing_{\mathcal F}$) et aux Inf pour $\leq$

On n'exige pas qu'elle transforme fig. élémentaire en figure élémentaire. Mais quand il en est ainsi, cette application est connue quand on connait sa restriction
\[
f_{\unsure{\mathrm{str}1|\mathrm{str}l}}:M'\unsure{\mathrel{\relbar\joinrel\relbar\joinrel\shortmid}|\longrightarrow}M
\]
aux ens. de \unsure{multistrates} \inserted{(= fig. élém.)}. Dans ce cas, pour que $f$ soit \unsure{injective|injection}, il f. et s. que $f_{\mathrm{str}1}^{\unsure{F|\mathcal F}}$ le soit (et en termes des plongements $\mathcal F\hookrightarrow\unsure{\mathcal P|\mathrm P}(M)$, $\mathcal F'\unsure{\lhook\joinrel\relbar|\hookrightarrow}\mathcal P(M')$, $f$ s'identifie alors à la restriction \deleted{\unsure{/corestriction}} de \deleted{$\unsure{f}$} $A\to f_{\mathrm{str}1}(A)$ : $\mathcal P(M')\longrightarrow\mathcal P(M)$, du moins quand $f(M')$ est une partie \inserted{$\leq$-fermée de $M$, cf plus bas.}
% slanted
\margin{pas clair, il \deleted{\unsure{pour}} faudrait que \inserted{$F'\in\mathcal F'$}, on ait $\widetilde{f(F')}=f(\widetilde{F'})$, ce qui n'est pas clair, cf plus bas.}

Il faudrait arriver à imposer aux quasi-morph. des conditions \unsure{supplt|suppl}, qui impliquent que lorsque $f$ est injectif, alors l'application induite :
